\chapter{The Rig in Practice: A Claude Code Reference}
\label{chapter:RigReference}
% EVOLVE-BLOCK-START

This appendix is a working reference for Claude Code, the agentic CLI used throughout the course --- the concrete counterpart to the rig concepts of Chapter~\ref{chapter:AgenticRig}.\index{Claude Code}
The chapters deliberately teach concepts that outlive any particular tool; this appendix deliberately does the opposite, collecting the command-level detail you need at the keyboard.
Product surfaces change faster than concepts: treat this appendix as a snapshot (Fall 2026), verify against the current documentation at \href{https://docs.anthropic.com}{docs.anthropic.com} when something disagrees, and see the course's vendored workshop (\code{resources/claude-cli-workshop/}) for the guided version of this material with exercises.

\section{Launching and Sessions}

Claude Code runs \emph{in} your terminal and inherits your shell environment --- your \code{PATH}, your environment variables, your working directory, your git configuration.
The directory you launch from is the project context (Section~\ref{section:InstructionFiles}), so start it from the project root.

\begin{center}
\begin{tabular}{ll}
\toprule
Command & Effect \\
\midrule
\code{claude}                          & start an interactive session in this directory \\
\code{claude -c}                       & continue the most recent session here \\
\code{claude -r}                       & open the session picker (search, preview, resume) \\
\code{claude -p "..."}                 & non-interactive: run one prompt, print, exit \\
\code{claude -{}-model <name>}           & choose the model at launch \\
\code{claude -{}-permission-mode plan}   & start in read-only plan mode \\
\code{claude -{}-add-dir <path>}         & grant access to a directory outside the project \\
\bottomrule
\end{tabular}
\end{center}

A session's conversation history persists on disk and is restored on resume; the context-window cost of that history persists too, so long resumed sessions may need compaction (\code{/compact}, below).
Sessions are filed under the directory they were started from --- work from the project root and the session list stays organized by project.

\section{Permission Modes}

Chapter~\ref{chapter:AgenticRig} introduced the layered permission model (Principle~\ref{principle:LeastPrivilege}); Claude Code exposes it as three interactive modes, cycled with \code{Shift+Tab}:

\begin{itemize}
\item \textbf{Normal} --- the agent asks before editing files or running commands; you approve each action.
\item \textbf{Auto-accept} --- file edits apply without prompting; commands still follow the configured permission rules. The productive middle ground once a plan is agreed.
\item \textbf{Plan} --- read-only. The agent can explore and propose but cannot change anything. Use it for anything touching more than a couple of files, before granting write access.
\end{itemize}

The plan-then-execute workflow is the practical form of the specify--execute--read loop of Chapter~\ref{chapter:AgenticCoding}: enter plan mode, describe the goal, refine the proposed plan, then switch modes and let the agent execute the plan it already committed to.

Standing rules live in \code{.claude/settings.json}, which is committed to the repository so a team shares one trust configuration:

\begin{lstlisting}[language=json]
{
  "permissions": {
    "allow": ["Bash(git status)", "Bash(pytest *)"],
    "deny":  ["Read(.env)", "Read(secrets/**)"]
  }
}
\end{lstlisting}

The \code{allow} list removes prompts for named commands; the \code{deny} list makes files unreadable to the agent regardless of mode --- the right home for credentials and raw data that no session should touch.

\begin{remark}[The permission ladder has a middle rung]
Between Normal (approve every action) and a blanket skip of all prompts sits a graduated middle: OS-level sandboxing that bounds filesystem and network access, and --- as of 2026 --- an \emph{auto} mode whose action classifier auto-approves demonstrably safe commands and escalates only risky ones.
Both exist because approval is a finite resource: when nearly every prompt is approved by reflex, human-in-the-loop stops being a safeguard.
Treat the specific tiers as a fast-moving product surface and verify against current documentation; the durable idea is that autonomy and safety are set by \emph{where} the boundary is drawn, not by how many prompts a tiring human answers.
\end{remark}

\section{Instruction File Layers}

The instruction hierarchy of Section~\ref{section:InstructionFiles} is realized as stacked \code{CLAUDE.md} files; more specific layers extend or override broader ones.

\begin{center}
\begin{tabular}{llll}
\toprule
Layer & Location & Scope & Shared? \\
\midrule
Global       & \code{\textasciitilde/.claude/CLAUDE.md} & every project      & no (personal) \\
Project      & \code{./CLAUDE.md}                       & this repository    & yes (committed) \\
Local        & \code{./CLAUDE.local.md}                 & this repository    & no (gitignored) \\
Subdirectory & \code{src/api/CLAUDE.md}                 & that module        & yes (committed) \\
\bottomrule
\end{tabular}
\end{center}

Subdirectory files load on demand --- only when the agent works with files in that directory --- which keeps module-specific rules from occupying the window (Chapter~\ref{chapter:ContextEngineering}) on unrelated tasks.
\code{/init} generates a starter project file by analyzing the repository; edit it rather than accepting it verbatim.

\section{In-Session Controls}

Slash commands control the session itself rather than instructing the agent:

\begin{center}
\begin{tabular}{ll}
\toprule
Command & Effect \\
\midrule
\code{/init}     & generate a starter \code{CLAUDE.md} for the project \\
\code{/model}    & switch models mid-session \\
\code{/context}  & show how much of the context window is in use \\
\code{/compact}  & summarize the conversation to reclaim window space \\
\code{/clear}    & reset the conversation entirely \\
\code{/cost}     & show token usage and cost for the session \\
\code{/help}     & list available commands \\
\bottomrule
\end{tabular}
\end{center}

Two prefixes are worth internalizing.
\code{!}~runs a shell command directly, without the agent interpreting it (\code{! git status}) --- useful for quick checks mid-session.
\code{@}~pulls a file into focus (\code{@src/eval.py what does the scorer return?}), a manual form of the retrieval this course automates in Chapter~\ref{chapter:Retrieval}.

\code{/context} and \code{/compact} are the manual face of the context management developed in Chapter~\ref{chapter:ContextEngineering}: check what occupies the window, and compact --- accepting the summarization loss --- when it fills.
Starting a fresh session on a task switch is often the better move than compacting a sprawling one.

\section{Headless Operation}

The \code{-p} flag runs the agent non-interactively --- the building block for the scripted, unattended pipelines of Chapter~\ref{chapter:Orchestration}:

\begin{lstlisting}[language=bash]
claude -p "Fix the failing tests, then stop." \
  --max-turns 10 \
  --allowedTools "Read" "Edit" "Bash(pytest *)" \
  --output-format json > result.json
\end{lstlisting}

Three flags do the safety work.
\code{-{}-max-turns} bounds the loop (the timeout discipline of Chapter~\ref{chapter:Orchestration}); \code{-{}-allowedTools} grants exactly the tools the task needs (least privilege, made per-invocation); \code{-{}-output-format json} makes the result machine-readable so a harness --- not a human --- can consume it.
An unattended agent has no human in the loop, so the permission configuration \emph{is} the supervision: headless runs deserve narrower allowlists than interactive ones, not broader.

\section{Parallel Work with Git Worktrees}

Git \emph{worktrees}\index{Worktree} give one repository several simultaneous working directories, each with its own checked-out branch and shared history:

\begin{lstlisting}[language=bash]
git worktree add ../proj-tests feature/add-tests
git worktree add ../proj-docs  feature/update-docs
# one agent per worktree, in separate terminals:
#   (cd ../proj-tests && claude -p "Add the missing unit tests")
#   (cd ../proj-docs  && claude -p "Update the API documentation")
git merge feature/add-tests   # then merge and remove the worktrees
git worktree remove ../proj-tests
\end{lstlisting}

This is the physical infrastructure for the multi-agent patterns of Chapter~\ref{chapter:Orchestration}: each parallel agent gets an isolated directory, so agents cannot overwrite each other's changes, and every result arrives as a branch a human can review before merging.

\section{Extending the Agent}

Three extension mechanisms recur in the course, listed here with their conceptual homes:

\begin{itemize}
\item \textbf{MCP servers} connect the agent to systems beyond the filesystem --- databases, issue trackers, web services --- as additional tools, using the protocol introduced in Chapter~\ref{chapter:ContextEngineering}. Configure per-project so the connection is shared and version-controlled; every added tool is also added attack surface for the prompt-injection risks of Chapter~\ref{chapter:AgenticRig}.
\item \textbf{Skills} are markdown procedure files the agent reads before performing a recurring task --- a checklist it follows, invoked as a slash command. They are the lightweight precursor to the tool definitions of Chapter~\ref{chapter:AgentLoop}: instructions, not code.
\item \textbf{Hooks} run your scripts at fixed points in the agent's lifecycle (before or after tool calls), enforcing policy mechanically --- lint after every write, log every command for audit. They implement guardrails the model cannot forget, because the harness, not the model, executes them.
\end{itemize}

Configuration details for all three change frequently; consult the current documentation rather than memorizing syntax.
Skills in particular lean on \emph{progressive disclosure}: a skill preloads only its name and description, and loads its full \code{SKILL.md} body and bundled files only when the agent judges the task relevant, so the extension surface does not tax the window (Chapter~\ref{chapter:ContextEngineering}) until it is used (Anthropic, 2025).
The durable idea is the division of labor: instruction files shape \emph{behavior}, permissions bound \emph{authority}, and hooks enforce \emph{invariants}.

\section*{Further Reading}

\begin{small}
\begin{enumerate}
\item Anthropic, ``Claude Code: Best Practices for Agentic Coding,'' \emph{Anthropic Engineering}, 2025 --- the practitioner reference for the command-level surface snapshotted here; a moving product surface, now at \href{https://code.claude.com/docs}{code.claude.com/docs}.
\item Anthropic, ``How We Built Claude Code Auto Mode: A Safer Way to Skip Permissions,'' \emph{Anthropic Engineering}, 2026 --- the \emph{auto} permission tier: a classifier that auto-approves demonstrably safe actions and escalates risky ones.
\item Anthropic, ``Beyond Permission Prompts: Making Claude Code More Secure and Autonomous,'' \emph{Anthropic Engineering}, 2025 --- OS-level filesystem and network sandboxing as the structural counterpart to per-action permissions.
\item Anthropic, ``Equipping Agents for the Real World with Agent Skills,'' \emph{Anthropic Engineering}, 2025 --- Skills as filesystem folders (\code{SKILL.md} plus bundled files) discovered by metadata and loaded through progressive disclosure.
\end{enumerate}
\end{small}

% EVOLVE-BLOCK-END
